%% ma: L11-B21-0005
%% lop: 11
%% chuong: 6
%% bai: 21
%% dang: 11.21.1
%% loai: TN4
%% muc: TH
%% dap_an: "B"
%% nguon: "(NBV)-ĐỀ SỐ 7-MỖI NGÀY 1 ĐỀ THI 2025.pdf (D07, MỖI NGÀY 1 ĐỀ - Đề số 7), Phần I Câu 7"
%% kien_thuc: "Bài 21 (bất phương trình lôgarit)"
%% trang_thai: chinh-thuc
%% ghi_chu: "Trùng ý tưởng với dạng 11.21.1 (bất phương trình lôgarit cơ bản; ở đây so sánh hai lôgarit cùng cơ số). Đáp án gốc B đúng."
\begin{ex}
Tập nghiệm của bất phương trình $\log_2(x+2)\ge\log_2(6-x)$ là
\choice{$(-2;2]$}{\True $[2;6)$}{$[2;+\infty)$}{$(-\infty;2]$}
\loigiai{Điều kiện $x+2>0$ và $6-x>0$, tức $-2<x<6$. Cơ số $2>1$ nên bất phương trình tương đương $x+2\ge6-x\Leftrightarrow x\ge2$. Kết hợp điều kiện: $[2;6)$.}
\end{ex}
